\documentclass[11pt,a4paper]{article}

\usepackage[utf8]{inputenc}
\usepackage[T1]{fontenc}
\usepackage[french]{babel}
\usepackage[margin=2.3cm]{geometry}
\usepackage{amsmath,amssymb,amsthm}
\usepackage{booktabs}
\usepackage{graphicx}
\usepackage{enumitem}
\usepackage[hidelinks]{hyperref}
\usepackage{setspace}
\usepackage{caption}
\usepackage{tikz}
\onehalfspacing

\newtheorem{theorem}{Théorème}[section]
\newtheorem{proposition}[theorem]{Proposition}
\newtheorem{lemma}[theorem]{Lemme}
\theoremstyle{remark}
\newtheorem{remark}[theorem]{Remarque}
\theoremstyle{definition}
\newtheorem{definition}[theorem]{Définition}

\DeclareMathOperator{\prox}{prox}
\newcommand{\R}{\mathbb{R}}
\newcommand{\pinball}{\rho_\tau}
\newcommand{\T}{^{\top}}

\begin{document}

%% ------------------------------------------------------------------ PAGE DE GARDE
\begin{titlepage}
\thispagestyle{empty}
\begin{tikzpicture}[remember picture,overlay]
  \node[anchor=north west,inner sep=0] at (current page.north west)
    {\includegraphics[width=\paperwidth,height=\paperheight]{fond_garde.png}};
  % En-tête (sous le logo ESPRIM)
  \node[anchor=north,align=center,text width=15cm] at ([yshift=-9.9cm,xshift=10.5cm]current page.north west)
    {\large\bfseries Rapport de projet\\[2pt]
     \normalsize Optimisation pour le Machine Learning --- 5e année Data Sciences et IA};
  % Titre du projet
  \node[anchor=north,align=center,text width=16cm] at ([yshift=-11.6cm,xshift=10.5cm]current page.north west)
    {\Large\bfseries Prévision probabiliste de l'énergie solaire en Afrique\\[5pt]
     \large Régression quantile sous contraintes physiques};
  % Année universitaire
  \node[anchor=north,align=center] at ([yshift=-14.6cm,xshift=10.5cm]current page.north west)
    {\bfseries Année universitaire 2026 / 2027};
  % Auteurs et encadrant
  \node[anchor=north,align=center,text width=16cm] at ([yshift=-17.5cm,xshift=10.5cm]current page.north west)
    {\textbf{Réalisé par :}\\[2pt]
     Kawthar Jguirim, Ons Bel Hej Saleh, Mouheb Sayadi,\\
     Yosr Bhouri, Mohamed Abidi};
  \node[anchor=north,align=center] at ([yshift=-20.3cm,xshift=10.5cm]current page.north west)
    {\textbf{Encadrant :} Bilel Charfi};
  \node[anchor=north,align=center] at ([yshift=-22.4cm,xshift=10.5cm]current page.north west)
    {\small Module : Optimisation pour le Machine Learning\\[2pt]
     \small Date de soutenance : \textit{à compléter}};
  \node[anchor=south,align=center] at ([yshift=2.2cm,xshift=10.5cm]current page.south west)
    {\small École Supérieure Privée d'Ingénieurs de Monastir};
\end{tikzpicture}
\end{titlepage}

%% ------------------------------------------------------------------ RESUME
\section*{Résumé}
Ce rapport présente la prévision probabiliste du rayonnement solaire journalier (\texttt{ALLSKY\_SFC\_SW\_DWN}, en kWh/m$^2$/j)
pour le lendemain, sur trois sites d'Afrique du Nord et de l'Ouest (Monastir, Tamanrasset, Dakar). La prévision est fournie
sous forme de trois quantiles ($\tau = 0{,}1$ ; $0{,}5$ ; $0{,}9$), à partir de données NASA POWER de 2014 à juin 2026.
Quatre axes sont traités : (1) régression ridge et descente à pas fixe, pas exact, Armijo, Wolfe et gradient conjugué ;
(2) régression quantile non dérivable, par sous-gradient projeté, lissage de Moreau et Newton amorti ;
(3) optimisation sous contraintes de non-croisement et de bornes de ciel clair, par dualité et lagrangien augmenté ;
(4) piste R3 : prédiction conforme (CQR et ACI). Les algorithmes sont écrits à la main en NumPy ; scikit-learn, SciPy et cvxpy
ne servent qu'à la vérification.

Principaux résultats : le ridge bat la prévision naïve sur les trois sites (RMSE réduite de 12 à 22 \%) ; le lagrangien
augmenté impose les contraintes en 5 itérations externes ; la régression quantile sans contrainte dépasse le ciel clair sur 71 \% des jours
d'apprentissage à Monastir, ce qui la rend physiquement incohérente ; CQR porte la couverture de l'intervalle à 80 \% à 81,1 \% à Monastir,
mais seulement à 76,7 \% à Tamanrasset, où la calibration ne reflète pas bien le test.

%% ------------------------------------------------------------------ INTRO
\clearpage
\tableofcontents
\clearpage

%% ------------------------------------------------------------------ INTRODUCTION GENERALE
\section{Introduction générale}
\subsection{Contexte et problème}
Une centrale solaire produit d'autant plus d'électricité que le rayonnement reçu au sol est fort. L'opérateur du réseau doit
décider la veille du niveau de production attendu : s'il surestime le rayonnement, il manque d'électricité ; s'il le sous-estime,
il brûle du combustible inutilement. Une prévision ponctuelle ne suffit donc pas : il faut une fourchette.

Nous prévoyons $Y_t = \mathrm{ALLSKY}(t+1)$ à partir des variables connues le jour $t$ (température, humidité, vent, pression,
précipitations, nébulosité, trois jours de rayonnement passé et saisonnalité). La borne physique $c_t = \mathrm{CLRSKY}(t+1)$, rayonnement
par ciel clair calculé astronomiquement, sert de contrainte supérieure ; elle n'est jamais une variable explicative.

\subsection{Objectifs}
Les objectifs fixés par le cahier sont : formuler chaque problème et dire s'il admet une solution unique ; choisir un algorithme
et ses réglages grâce à la théorie, puis vérifier la vitesse observée ; livrer une prévision probabiliste par site ; comparer les sites
et argumenter les conclusions.

\subsection{Organisation du rapport}
La section~\ref{sec:etat} rappelle l'état de l'art. La section~\ref{sec:donnees} décrit les données et les règles d'absence de fuite.
Les sections~\ref{sec:ax1}, \ref{sec:ax2} et \ref{sec:ax3} traitent les axes 1 à 3 avec leurs démonstrations, la section~\ref{sec:ax4}
la piste R3, et la section~\ref{sec:disc} discute les résultats et leurs limites.

%% ------------------------------------------------------------------ ETAT DE L'ART
\section{État de l'art}\label{sec:etat}
\paragraph{Régression quantile.} Koenker et Bassett (1978) introduisent la régression quantile : minimiser la perte pinball
$\pinball(r) = \max(\tau r, (\tau-1)r)$ donne un estimateur du $\tau$-quantile conditionnel. Cette perte est convexe mais non dérivable.
Koenker (2005) en donne la théorie et les algorithmes de programmation linéaire.

\paragraph{Lissage et sous-gradient.} Le lissage de Moreau remplace une fonction convexe par son enveloppe, dont le gradient est
lipschitzien. Shor (1985) et Nemirovski et Yudin (1983) analysent la méthode du sous-gradient projeté ; Lacoste-Julien, Schmidt et Bach (2012)
donnent une vitesse $O(1/k)$ avec moyenne pondérée des itérés pour les fonctions fortement convexes.

\paragraph{Optimisation sous contraintes.} Les conditions de Karush, Kuhn et Tucker (KKT) caractérisent les solutions. Le dual
d'un problème fortement convexe est lisse ; la montée de gradient projeté sur le dual converge en $O(1/k)$, et le lagrangien augmenté
(Hestenes 1969 ; Powell 1969 ; Rockafellar 1973) évite la mauvaise conditionnement du dual.

\paragraph{Prévision conforme.} Vovk et al. (2005) proposent la prédiction conforme ; Romano, Patterson et Candès (2019) introduisent CQR,
qui combine régression quantile et calibration. Gibbs et Candès (2021) proposent ACI, une mise à jour en ligne du niveau de couverture avec une borne
sur la couverture moyenne.

\paragraph{Données et climat.} NASA POWER fournit des données journalières de rayonnement et de météorologie sur grille, sans compte.
La classification de Köppen-Geiger est recalculée selon Beck et al. (2018) à partir des moyennes mensuelles 1991--2020.

%% ------------------------------------------------------------------ DONNEES
\section{Données et méthodologie}\label{sec:donnees}
\subsection{Sites et période}
Trois sites sont retenus, pour contraste climatique : Monastir (Tunisie, 35,78\textdegree{}N ; 10,83\textdegree{}E), Tamanrasset (Algérie, 22,79\textdegree{}N ; 5,52\textdegree{}E),
et Dakar (Sénégal, 14,69\textdegree{}N ; $-17{,}44$\textdegree{}E). Les données couvrent 2014-01-01 à 2026-06-30, soit 4\,564 jours sans lacune ni doublon.
L'apprentissage porte sur les cibles de 2014 à 2023 (3\,649 jours), le test sur les cibles de 2024 à juin 2026 (912 jours).

\subsection{Classification climatique}
La classe de Köppen-Geiger de chaque site est donnée par la carte de Beck et al. (2018), via kgcpy : Monastir BSh, Tamanrasset BWh, Dakar BSh.
Le recalcul à partir des moyennes NASA 1991--2020 (critères de Beck et al.) donne les mêmes classes, sans écart à expliquer.

\subsection{Variables et absence de fuite}
Les entrées du jour $t$ sont : $T_2M$, $T_2M_{\max}$, $T_2M_{\min}$, $RH_{2M}$, $WS_{2M}$, $PS$, $PRECTOTCORR$, $CLOUD\_AMT$,
trois valeurs de rayonnement ($t$, $t-1$, $t-2$) et deux termes de saisonnalité ($\sin$, $\cos$ du jour de l'année).
La cible est $Y_t = \mathrm{ALLSKY}(t+1)$. Les règles suivantes évitent toute information future :
\begin{itemize}[itemsep=2pt]
\item la séparation apprentissage/test se fait sur la date de la cible : la ligne du 31/12/2023 prévoit le 01/01/2024 et appartient donc au test ;
\item le centrage de la cible, la standardisation et l'imputation des valeurs manquantes (propagation vers l'avant) n'utilisent que l'apprentissage ;
\item $c_t = \mathrm{CLRSKY}(t+1)$ n'est utilisée que comme borne.
\end{itemize}
La matrice de conception $A \in \R^{n\times d}$ contient une colonne de 1 ajoutée après standardisation, donc $d = 14$ pour les axes 1 à 3.

%% ------------------------------------------------------------------ AXE 1
\section{Axe 1 : régression ridge}\label{sec:ax1}
\subsection{Problème et constantes}
On minimise
\begin{equation}\label{eq:ridge}
f(w) = \frac{1}{2n}\|Aw - y\|^2 + \frac{\lambda}{2}\|w\|^2, \qquad \lambda = 10^{-2}.
\end{equation}
Le gradient est $\nabla f(w) = \frac1n A\T(Aw-y) + \lambda w$ et la hessienne $H = \frac1n A\T A + \lambda I$.

\begin{proposition}[Existence, unicité, constantes]\label{prop:ridge}
Soit $\mu = \lambda_{\min}(H)$ et $L = \lambda_{\max}(H)$. Alors $f$ est $\mu$-fortement convexe, $\nabla f$ est $L$-lipschitzien,
$f$ admet un unique minimiseur $w^\star = H^{-1} A\T y / n$, et $\kappa = L/\mu$.
En particulier, $\mu \ge \lambda > 0$ pour $\lambda > 0$, et pour $\lambda = 0$, $\mu = \lambda_{\min}(A\T A / n) > 0$ dès que $A$ est de rang plein.
\end{proposition}
\begin{proof}
La hessienne est $H \succeq \lambda I \succ 0$, donc $f$ est fortement convexe de module $\mu$, et $\nabla f$ est $L$-lipschitzien car $\|H\|_2 = L$.
Le minimiseur vérifie $\nabla f(w^\star) = 0$, soit $H w^\star = A\T y / n$, d'où l'expression. L'unicité découle de la stricte convexité.
\end{proof}

\subsection{Descente à pas fixe et vitesse}
\begin{proposition}[Vitesse à pas fixe]\label{prop:fixe}
Pour $w_{k+1} = w_k - \alpha\nabla f(w_k)$ avec $0 < \alpha < 2/L$, on a
$w_{k+1} - w^\star = (I - \alpha H)(w_k - w^\star)$ et
\[
f(w_k) - f^\star \le \rho(\alpha)^{2k}\,\bigl(f(w_0) - f^\star\bigr), \qquad \rho(\alpha) = \max\bigl(|1-\alpha\mu|,\,|1-\alpha L|\bigr).
\]
Le pas optimal est $\alpha = 2/(L+\mu)$, pour lequel $\rho = (\kappa-1)/(\kappa+1)$. Le pas $\alpha = 1/L$ donne $\rho = 1-1/\kappa$.
\end{proposition}
\begin{proof}
Le gradient vaut $\nabla f(w) = H(w - w^\star)$, donc $w_{k+1} - w^\star = (I-\alpha H)(w_k-w^\star)$. Dans une base propre de $H$,
$f(w) - f^\star = \frac12 (w-w^\star)\T H (w-w^\star)$, et chaque composante propre $\lambda_i \in [\mu, L]$ est multipliée par $(1-\alpha\lambda_i)^{2}$ après $k$ itérations au carré.
Le maximum sur $\lambda_i$ de $|1-\alpha\lambda_i|$ est $\rho(\alpha)$, d'où la borne. Minimiser $\rho(\alpha)$ en $\alpha$ donne $\alpha = 2/(L+\mu)$.
\end{proof}

\begin{remark}
Le nombre d'itérations pour diviser l'écart par $10^6$ est $k \approx \ln(10^6)/(-2\ln\rho)$, soit environ $6{,}9\,\kappa$ pour $\alpha = 1/L$
et $3{,}5\,\kappa$ pour $\alpha = 2/(L+\mu)$. C'est une borne au pire cas : elle est atteinte seulement si l'erreur initiale a une composante
dans la direction de courbure $\mu$.
\end{remark}

\subsection{Pas exact, Armijo, Wolfe}
\begin{proposition}[Pas exact et Armijo]\label{prop:exact}
Le pas exact $\alpha_k = \|g_k\|^2 / (g_k\T H g_k)$ satisfait la condition d'Armijo $f(w_k - \alpha_k g_k) \le f(w_k) - c_1\alpha_k\|g_k\|^2$
pour tout $c_1 \le 1/2$.
\end{proposition}
\begin{proof}
Sur une quadratique, $\varphi(\alpha) = f(w - \alpha g) = f(w) - \alpha\|g\|^2 + \frac{\alpha^2}{2} g\T H g$. En $\alpha_k$,
$\varphi(\alpha_k) = f(w) - \frac{\|g\|^4}{2\,g\T H g}$. Or $\alpha_k\|g\|^2 = \|g\|^4/(g\T H g)$, donc $\varphi(\alpha_k) = f(w) - \tfrac12\alpha_k\|g\|^2$,
ce qui est $\le f(w) - c_1\alpha_k\|g\|^2$ dès que $c_1 \le 1/2$.
\end{proof}
Les recherches d'Armijo (rétrécissement de $1$, $c_1 = 10^{-4}$) et de Wolfe forte ($c_1 = 10^{-4}$, $c_2 = 0{,}9$, par encadrement puis dichotomie)
sont implémentées dans le notebook et comptées en évaluations de $f$ et de gradient.

\subsection{Gradient conjugué}
Le gradient conjugué résout $Hw = A\T y/n$ et atteint, après $k$ itérations, une erreur d'énergie
$\|e_k\|_H \le 2\bigl(\tfrac{\sqrt\kappa-1}{\sqrt\kappa+1}\bigr)^k\|e_0\|_H$. Pour un écart de $10^{-6}$ sur $f$ (soit $10^{-3}$ sur $\|e\|_H$), le nombre
d'itérations est donc de l'ordre de $\tfrac12\sqrt\kappa\,\ln(2\cdot10^3)$.

\subsection{Équivalence avec scikit-learn}
\begin{proposition}\label{prop:sk}
Soit $\alpha_{sk} = n\lambda$ et \texttt{fit\_intercept=False} dans \texttt{sklearn.linear\_model.Ridge}. Alors la solution de \texttt{Ridge} est $w^\star$.
\end{proposition}
\begin{proof}
\texttt{Ridge} minimise $\|y - Aw\|^2 + \alpha_{sk}\|w\|^2$. Multiplier $f$ par $2n$ donne $\|Aw-y\|^2 + n\lambda\|w\|^2$.
Les deux problèmes ont donc les mêmes minimiseurs si $\alpha_{sk} = n\lambda$. L'intercept est déjà la colonne de 1 de $A$, d'où \texttt{fit\_intercept=False}.
\end{proof}

\subsection{Résultats (Monastir, $\lambda = 10^{-2}$)}
\begin{table}[h]\centering
\begin{tabular}{lrrr}
\toprule
Méthode & Itérations prédites & Itérations observées & Évaluations \\
\midrule
Pas fixe $1/L$ & 3\,895 & 858 & 0 \\
Pas fixe $2/(L+\mu)$ & 1\,950 & 1\,915 & 0 \\
Pas exact & 1\,950 & 273 & 0 \\
Armijo & -- & 391 & 643 \\
Wolfe & -- & 391 & 1\,260 \\
Gradient conjugué & 91 & 14 & 0 \\
\bottomrule
\end{tabular}
\caption{Itérations pour diviser $f - f^\star$ par $10^6$. Les prédictions sont des bornes au pire cas.}
\end{table}
Pour $\lambda = 10^{-2}$, $\kappa = 564{,}3$ ($L = 6{,}402$, $\mu = 0{,}01134$). L'écart entre le ridge de la proposition~\ref{prop:sk} et \texttt{Ridge} vaut $6\cdot10^{-15}$
sur les coefficients. Le préconditionnement diagonal ne change pas $\kappa$ (la diagonale de $H$ vaut 1{,}01 partout) : pas de gain.
La prévision ridge atteint une RMSE de 0,939 kWh/m$^2$/j contre 1,204 pour la prévision naïve.

%% ------------------------------------------------------------------ AXE 2
\section{Axe 2 : régression quantile}\label{sec:ax2}
\subsection{Problème, sous-différentiel, unicité}
Pour $\tau \in \{0{,}1 ; 0{,}5 ; 0{,}9\}$, on minimise
\begin{equation}\label{eq:quant}
F(w) = \frac1n\sum_{t=1}^n \pinball(y_t - a_t\T w) + \frac\lambda2\|w\|^2.
\end{equation}
Le sous-différentiel de $\pinball$ est $\{\tau\}$ si $r>0$, $\{\tau-1\}$ si $r<0$, et $[\tau-1,\tau]$ si $r=0$.
Pour $\psi_t \in \partial\pinball(r_t)$, un sous-gradient de $F$ est $g(w) = -\frac1n\sum_t \psi_t a_t + \lambda w$.

\begin{proposition}\label{prop:quant}
$F$ est $\lambda$-fortement convexe et coercive. Elle admet un unique minimiseur $w^\star$, et $\|w^\star\| \le R := \sqrt{2F(0)/\lambda}$.
\end{proposition}
\begin{proof}
$\pinball$ est convexe, donc $F$ est la somme d'une fonction convexe et de $\frac\lambda2\|w\|^2$, donc $\lambda$-fortement convexe.
Comme $F(w^\star) \le F(0)$ et $\frac\lambda2\|w^\star\|^2 \le F(w^\star) - \min_w \frac1n\sum\pinball \le F(0)$, on obtient la borne sur $\|w^\star\|$.
\end{proof}

\subsection{Sous-gradient projeté}
Sur la boule de rayon $R$, $\|g\| \le G := \max_t\|a_t\|\cdot\max(\tau,1-\tau) + \lambda R$. L'algorithme est
\[
w_{t+1} = \Pi_R\bigl(w_t - \eta_t g_t\bigr), \qquad \eta_t = \frac{2}{\lambda(t+1)}, \qquad \bar w_T = \frac{\sum_{t=1}^T t\,w_t}{\sum_{t=1}^T t}.
\]
\begin{theorem}\label{thm:subgrad}
Pour $\mu = \lambda$ et $\eta_t = 2/(\mu(t+1))$, on a $F(\bar w_T) - F^\star \le \dfrac{2G^2}{\mu(T+1)}$.
\end{theorem}
\begin{proof}[Esquisse]
C'est la borne de Lacoste-Julien, Schmidt et Bach (2012), appliquée à une fonction $\mu$-fortement convexe dont les sous-gradients sont bornés par $G$ sur l'ensemble admissible.
La moyenne pondérée par $t$ est celle qui donne le terme en $1/T$. Nous ne reproduisons pas la preuve complète.
\end{proof}

\subsection{Lissage de Moreau}
Soit $\varepsilon > 0$. L'enveloppe de Moreau de $\pinball$ est
\[
M_\varepsilon(r) = \begin{cases}\tau r - \varepsilon\tau^2/2, & r > \varepsilon\tau,\\[2pt] r^2/(2\varepsilon), & \varepsilon(\tau-1) \le r \le \varepsilon\tau,\\[2pt] (\tau-1) r - \varepsilon(1-\tau)^2/2, & r < \varepsilon(\tau-1).\end{cases}
\]
\begin{proposition}\label{prop:moreau}
$M_\varepsilon$ est convexe, de dérivée $M'_\varepsilon(r) \in [\tau-1,\tau]$ $1/\varepsilon$-lipschitzienne, et
$0 \le \pinball(r) - M_\varepsilon(r) \le \varepsilon\max(\tau,1-\tau)^2/2$.
\end{proposition}
\begin{proof}
$M_\varepsilon(r) = \min_u\{\pinball(u) + (r-u)^2/(2\varepsilon)\}$, donc la minimisation de $u$ donne $u = r - \varepsilon\tau$ si $r > \varepsilon\tau$,
$u = r + \varepsilon(1-\tau)$ si $r < \varepsilon(\tau-1)$, et $u = 0$ sinon, ce qui reproduit les trois cas. La dérivée est $\tau$, $r/\varepsilon$ ou $\tau-1$, d'où la lipschitzianité.
Pour l'écart, $\pinball - M_\varepsilon$ est constante égale à $\varepsilon\tau^2/2$ (ou $\varepsilon(1-\tau)^2/2$) dans chaque queue et nulle au centre, donc majorée.
\end{proof}
La fonction $F_\varepsilon(w) = \frac1n\sum M_\varepsilon(y_t - a_t\T w) + \frac\lambda2\|w\|^2$ est $\lambda$-fortement convexe, de gradient $L_\varepsilon$-lipschitzien avec
$L_\varepsilon = \lambda_{\max}(A\T A/n)/\varepsilon + \lambda$, et $|F - F_\varepsilon| \le \varepsilon\max(\tau,1-\tau)^2/2$.

\subsection{Newton amorti et identité d'intercept}
La hessienne de $F_\varepsilon$ vaut $H_\varepsilon(w) = \frac1{n\varepsilon}A\T D_\varepsilon(w) A + \lambda I \succeq \lambda I$, où $D_\varepsilon$ est la matrice diagonale
des indicateurs de $\{\varepsilon(\tau-1) \le r_t \le \varepsilon\tau\}$. Elle est donc toujours inversible : le pas de Newton amorti est défini,
avec convergence globale linéaire et locale plus rapide.

\begin{proposition}[Identité d'intercept]\label{prop:intercept}
Notons $w_0$ la composante d'intercept (dernière coordonnée). En $w^\star$, $\frac1n\#\{r_t<0\} \approx \tau - \lambda w_0^\star$, plus précisément
$\lambda w_0^\star \in \bigl[\tfrac1n(\tau n_+ + (\tau-1)(n_- + n_0)),\ \tfrac1n(\tau(n_+ + n_0) + (\tau-1)n_-)\bigr]$, où $n_+$, $n_-$, $n_0$ comptent les résidus positifs,
négatifs et proches de 0.
\end{proposition}
\begin{proof}
La condition $0 \in \partial F(w^\star)$ s'écrit $\frac1n\sum_t\psi_t a_t = \lambda w^\star$. La dernière composante de $a_t$ vaut 1, donc $\frac1n\sum_t\psi_t = \lambda w_0^\star$.
Les résidus positifs donnent $\psi_t = \tau$, les négatifs $\tau-1$, et les résidus nuls $\psi_t \in [\tau-1,\tau]$ librement.
\end{proof}

\subsection{Résultats (Monastir)}
\begin{table}[h]\centering
\begin{tabular}{lrrr}
\toprule
Quantile $\tau$ & $F^\star$ (cvxpy) & Écart Newton ($\varepsilon=10^{-3}$) & Part des résidus négatifs \\
\midrule
0,1 & 0,186095 & $4{,}3\cdot10^{-7}$ & 0,10962 \\
0,5 & 0,300670 & $1{,}6\cdot10^{-7}$ & 0,49767 \\
0,9 & 0,110733 & $4{,}4\cdot10^{-7}$ & 0,89230 \\
\bottomrule
\end{tabular}
\caption{Référence cvxpy (CLARABEL, tolérance $10^{-10}$) et précision de Newton amorti. La part de résidus négatifs vérifie
l'identité d'intercept à la précision du solveur.}
\end{table}
Le sous-gradient projeté atteint $F(\bar w) - F^\star = 1{,}4\cdot10^{-7}$ en $20\,000$ itérations, alors que la borne du théorème~\ref{thm:subgrad} vaut $1{,}57$.
La borne est très lâche : la moyenne pondérée suffit en pratique. Pour le lissage, la descente à pas fixe n'atteint pas le critère en $20\,000$ itérations
pour $\varepsilon \le 0{,}1$, alors que Newton amorti converge en 10 à 50 itérations.
À budget égal de $7{,}3$ millions de lignes, les mini-lots de taille 64 atteignent $4{,}3\cdot10^{-6}$, contre $5{,}6\cdot10^{-5}$ pour le déterministe.
Pour $\lambda = 0$, le programme linéaire donne la même valeur que \texttt{QuantileRegressor(alpha=0)} et que cvxpy, à $10^{-15}$ près.
Les trois quantiles ne se croisent jamais sur le test à Monastir ; la couverture de $q_{0,1}$ et $q_{0,9}$ vaut 12,3 \% et 91,0 \%.

%% ------------------------------------------------------------------ AXE 3
\section{Axe 3 : optimisation sous contraintes}\label{sec:ax3}
\subsection{Formulation}
On note $W = (w_{0,1}, w_{0,5}, w_{0,9}) \in \R^{3d}$ et
\[
\min_W \Phi(W) = \sum_\tau F_\tau(w_\tau) \quad \text{s.c.}\quad CW \le b,
\]
où les $8n$ lignes de $C$ sont : non-croisement $a_t\T(w_{0,1} - w_{0,5}) \le 0$ et $a_t\T(w_{0,5} - w_{0,9}) \le 0$ ; minimum $-a_t\T w_\tau \le \mu_y$ ;
maximum $a_t\T w_\tau \le c_t - \mu_y$, pour $\tau \in \{0,1 ; 0,5 ; 0,9\}$. Ici $\mu_y$ est la moyenne de la cible d'apprentissage.
La contrainte de borne traduit $0 \le q_\tau \le c_t$ pour les quantiles bruts $q_\tau = a_t\T w_\tau + \mu_y$.

\begin{proposition}[Existence et unicité]
Le problème admet une solution unique $W^\star$. Le point $W_0$ dont chaque $w_\tau$ n'a que l'intercept non nul, égal à $k_\tau - \mu_y$ avec
$k_\tau = (0{,}25\,;\,0{,}5\,;\,0{,}75)\cdot\min_t c_t$, est strictement admissible : la contrainte de non-croisement y est $k_{0,1} - k_{0,5} < 0$, les bornes $k_\tau > 0$ et $k_\tau < c_t$.
\end{proposition}
\begin{proof}
$\Phi$ est $\lambda$-fortement convexe et l'ensemble admissible est un polyèdre convexe fermé non vide ; donc le minimiseur existe et est unique.
Pour Monastir, $\min_t c_t = 2{,}052$ kWh/m$^2$/j, donc $W_0$ a une marge minimale de $0{,}513$ : la condition de Slater est satisfaite.
\end{proof}

\subsection{KKT et dualité}
\begin{proposition}[Conditions KKT]
$W^\star$ est optimal si et seulement s'il existe $\nu^\star \ge 0$ tel que $\nabla\Phi(W^\star) + C\T\nu^\star = 0$ (au sens du sous-différentiel),
$CW^\star \le b$, et $\nu^\star_i (CW^\star - b)_i = 0$ pour tout $i$. La dualité forte tient par Slater.
\end{proposition}
Le dual $g(\nu) = \min_W\{\Phi(W) + \nu\T(CW-b)\}$ est concave, et son gradient $\nabla g(\nu) = CW(\nu) - b$ est lipschitzien de constante
$L_d = \|C\|_2^2/\lambda$ : le primal est $\lambda$-fortement convexe.

\subsection{Algorithmes}
La montée de gradient projeté sur le dual utilise le pas $1/L_d$ ; la variante de Nesterov accélère la convergence. Le lagrangien augmenté
minimise $\Phi(W) + \frac1{2\rho}\sum_i\bigl[(\nu_i + \rho(CW-b)_i)_+^2 - \nu_i^2\bigr]$ en $W$ puis met à jour $\nu \leftarrow (\nu + \rho(CW-b))_+$.
Le paramètre $\rho$ est pris égal à $L_\varepsilon$, courbure du primal lissé ; ce choix est une règle de notre part, non imposée par le cahier.
La minimisation interne est faite par Newton amorti, avec hessienne généralisée $\rho C\T D C$ où $D$ indique les contraintes actives.

\subsection{Résultats (Monastir)}
\begin{table}[h]\centering
\begin{tabular}{lrr}
\toprule
Méthode & Itérations & Violation finale \\
\midrule
Dual projeté (pas $1/L_d$) & 20\,000 (plafond) & $2{,}6$ \\
Dual accéléré (Nesterov) & 20\,000 (plafond) & $4{,}5\cdot10^{-2}$ \\
Lagrangien augmenté & \textbf{5} & $1{,}8\cdot10^{-10}$ \\
\bottomrule
\end{tabular}
\caption{Convergence vers la faisabilité. Pour le dual, $\|C\|_2^2 = 1{,}17\cdot10^5$ donc $L_d = 1{,}17\cdot10^7$ ; la prédiction est de $2{,}5\cdot10^{12}$ itérations.}
\end{table}
Le lagrangien augmenté satisfait les KKT à $10^{-8}$ près et atteint $\Phi - \Phi^\star = 8\cdot10^{-8}$ (avec $\Phi^\star = 2{,}08724$).

\subsection{Effet des contraintes sur le test}
Sans contrainte, le quantile $0{,}9$ dépasse le ciel clair sur $71\%$ des jours d'apprentissage et sur $665$ jours de test.
Imposer les bornes et le non-croisement fait s'effondrer la prévision : $q_{0,5} = q_{0,9}$ en moyenne, pinball médian $0{,}74$ contre $0{,}32$,
et couverture de l'intervalle $4{,}5\%$ contre $79\%$. Le problème n'est pas algorithmique : les KKT sont vérifiées. Quelques jours à fort écart pèsent sur tout le plan linéaire.
Il s'agit d'une limite du modèle linéaire sur la cible brute, non d'un défaut d'optimisation. Le renforcement par marges
($\delta_{nc} = 0{,}01$, $\delta_{ub} = 0{,}05$) ne corrige pas ce point.

%% ------------------------------------------------------------------ AXE 4
\section{Axe 4 : piste R3, prédiction conforme}\label{sec:ax4}
\subsection{CQR}
Soit $\hat q_{0,1}$ et $\hat q_{0,9}$ ajustés sur un sous-ensemble d'apprentissage propre, et $E_i = \max(\hat q_{0,1}(x_i) - y_i,\ y_i - \hat q_{0,9}(x_i))$ sur une calibration de taille $m$.
On pose $Q$ le quantile d'ordre $\lceil (m+1)(1-\alpha)\rceil/m$ des $E_i$, et l'intervalle est $[\hat q_{0,1}-Q,\ \hat q_{0,9}+Q]$.

\begin{theorem}[Couverture de CQR, Romano et al. 2019]
Si $(X_i,Y_i)$ de la calibration et la nouvelle observation $(X,Y)$ sont échangeables, alors
$\Pr\bigl(Y \in [\hat q_{0,1}(X) - Q,\ \hat q_{0,9}(X) + Q]\bigr) \ge 1-\alpha$.
\end{theorem}
\begin{proof}[Esquisse]
Les scores $E_1,\dots,E_m,E$ sont échangeables, donc le rang de $E$ parmi eux est uniforme. L'événement $Y\in$ intervalle équivaut à $E \le Q$,
et la définition de $Q$ à l'ordre $\lceil(m+1)(1-\alpha)\rceil$ donne $\Pr(E \le Q) \ge 1-\alpha$.
\end{proof}

\subsection{ACI}
Gibbs et Candès (2021) mettent à jour $\alpha_{t+1} = \alpha_t + \gamma(\alpha - \mathrm{err}_t)$, avec $\mathrm{err}_t = \mathbf{1}\{Y_t \notin C_t\}$.
La couverture moyenne satisfait
\[
\Bigl|\frac1T\sum_{t=1}^T \mathrm{err}_t - \alpha\Bigr| \le \frac{\max(\alpha, 1-\alpha) + \gamma}{\gamma T}.
\]

\subsection{Résultats (Monastir, 5 graines)}
\begin{table}[h]\centering
\begin{tabular}{lrr}
\toprule
Méthode & Couverture (cible 80\,\%) & Largeur (kWh/m$^2$/j) \\
\midrule
Brut (quantiles 0,1 et 0,9) & 78,7\,\% & 1,85 \\
Gaussien ridge $\pm 1{,}28\sigma$ & 82,4\,\% & 2,18 \\
CQR, calibration 1 an & 81,1\,\% $\pm$ 0,6 & 1,97 $\pm$ 0,02 \\
ACI, $\gamma = 0{,}005$ & 80,0\,\% & 1,91 \\
\bottomrule
\end{tabular}
\caption{Intervalles à 80\,\% sur le test (912 jours). Le gaussien sous-couvre au printemps (71\,\% en MAM).}
\end{table}
Calibrer sur 6 mois, 1 an ou 2 ans donne une couverture moyenne de 81,2\,\%, 81,1\,\% et 80,6\,\%, avec un écart-type qui baisse de 0,74 à 0,34 point.

\subsection{Trois sites}
CQR couvre 81,1\,\% à Monastir, 76,7\,\% à Tamanrasset et 83,5\,\% à Dakar. La garantie de CQR suppose l'échangeabilité entre calibration (2023) et test (2024--2026).
À Tamanrasset, ce n'est pas vérifié : la couverture passe sous la cible avec un écart-type faible, ce qui indique un décalage systématique et non du bruit.
Une calibration glissante, par ACI, est plus sûre dans ce cas.

%% ------------------------------------------------------------------ RESULTATS & DISCUSSION
\section{Résultats comparés et discussion}\label{sec:disc}
\subsection{Classement des sites}\label{sec:classement}
Le classement de difficulté, selon la RMSE du ridge, est Tamanrasset (0,709), Dakar (0,790), puis Monastir (0,939). Le même ordre vaut pour la médiane pinball et pour la largeur de l'intervalle.
Il change pour la fiabilité des intervalles, où Tamanrasset est le plus mal calibré. Les largeurs sont exprimées en kWh/m$^2$/j : une largeur relative au niveau moyen serait plus juste pour comparer les sites.

\subsection{Difficulté d'optimisation et de prévision}
Dakar a le conditionnement le plus favorable ($\kappa = 398$) mais la descente à pas $1/L$ y est la plus lente (1\,319 itérations, contre 858 à Monastir).
Le rapport observé/prédit passe de 0,22 à 0,48. Le conditionnement seul ne prédit donc pas la vitesse d'un site à l'autre : la répartition du spectre compte aussi.

\subsection{Hypothèse D.1}
Une hypothèse sur les sites les plus difficiles devait être formulée avant les essais. \textit{À compléter par le groupe, avec la date de rédaction.} Sans cette hypothèse datée, les résultats ne peuvent pas être présentés comme une confirmation.

\subsection{Recommandation pour un nouveau site}
Pour un exploitant, la méthode proposée est :
\begin{enumerate}[itemsep=2pt]
\item une prévision ponctuelle par ridge ($\lambda = 10^{-2}$), qui bat la prévision naïve sur les trois sites ;
\item une régression quantile lissée, résolue par Newton amorti ;
\item des intervalles CQR avec au moins un an de calibration récente pour les climats variables (BSh), et ACI pour les climats désertiques (BWh) où la calibration peut dériver ;
\item l'imposition des bornes de ciel clair par lagrangien augmenté, mais non sur la valeur brute : nous recommandons de modéliser le ratio $k_t = \mathrm{ALLSKY}(t+1)/\mathrm{CLRSKY}(t+1)$.
\end{enumerate}
Ces conclusions reposent sur trois sites, dont deux de même type climatique, et une seule année de test. Elles doivent être confirmées avant généralisation.

\subsection{Limites}
\begin{itemize}[itemsep=2pt]
\item Trois sites seulement ; deux sont de type BSh.
\item Un seul test (912 jours, 2024--2026) ; pas de répétition sur plusieurs années.
\item Les vitesses prédites sont des bornes au pire cas ; l'écart observé n'est pas expliqué au-delà du conditionnement.
\item La contrainte de ciel clair sur la valeur brute dégrade fortement la prévision ; la piste du ratio n'est pas testée.
\item Le lissage de Moreau introduit une erreur de l'ordre de $\varepsilon/2$ ; elle n'est pas corrigée dans les métriques de test.
\item La garantie CQR n'est pas valable à Tamanrasset.
\end{itemize}

\subsection{Pistes d'approfondissement}
La piste R3 est traitée. Les pistes R1 (apprentissage fédéré) et R2 (robustesse de Sagawa et al.) pourraient utiliser les trois sites comme exploitants, avec les mêmes variables.

%% ------------------------------------------------------------------ CONCLUSION
\section{Conclusion}
Les quatre axes donnent une prévision probabiliste de l'irradiation journalière sur trois sites. Les algorithmes, écrits à la main, ont été vérifiés contre les
références (scikit-learn, SciPy, cvxpy). Le ridge et la régression quantile lissée sont fiables et rapides. Les contraintes physiques sur la valeur brute
ne sont pas adaptées au modèle linéaire, et la calibration CQR doit être contrôlée site par site.

%% ------------------------------------------------------------------ REFERENCES
\begin{thebibliography}{99}
\bibitem{koenker1978} R. Koenker, G. Bassett. \emph{Regression Quantiles}. Econometrica, 46(1), 1978.
\bibitem{koenker2005} R. Koenker. \emph{Quantile Regression}. Cambridge University Press, 2005.
\bibitem{lacoste2012} S. Lacoste-Julien, M. Schmidt, F. Bach. \emph{A simpler approach to obtaining an O(1/t) convergence rate for projected stochastic subgradient descent}. arXiv:1212.2002, 2012.
\bibitem{nemirovski1983} A. Nemirovski, D. Yudin. \emph{Problem Complexity and Method Efficiency in Optimization}. Wiley, 1983.
\bibitem{romano2019} Y. Romano, E. Patterson, E. Candès. \emph{Conformalized Quantile Regression}. NeurIPS, 2019.
\bibitem{gibbs2021} I. Gibbs, E. Candès. \emph{Adaptive Conformal Inference Under Distribution Shift}. NeurIPS, 2021.
\bibitem{vovk2005} V. Vovk, A. Gammerman, G. Shafer. \emph{Algorithmic Learning in a Random World}. Springer, 2005.
\bibitem{beck2018} H. E. Beck et al. \emph{Present and future Köppen-Geiger climate classification maps at 1-km resolution}. Scientific Data, 5, 2018.
\bibitem{sagawa2020} S. Sagawa, P. W. Koh, T. B. Hashimoto, P. Liang. \emph{Distributionally Robust Neural Networks}. ICLR, 2020.
\bibitem{nasapower} NASA POWER. \emph{Prediction Of Worldwide Energy Resources}. \url{https://power.larc.nasa.gov}.
\end{thebibliography}

\end{document}
